% GENERATED FILE. Edit canonical lessons, then run npm run generate:books.
% canonical-source-hash: fnv1a64:0f56e997e8078200
% canonical-lessons: ES-C16-ver

\chapter{Ver --- Seeing}
\label{ch:spanish-ver}
\hlchaptermodality{53}

\begin{chapteropening}
\textbf{By the end of this chapter:} \emph{I can say what I see, with a verb whose stem keeps a letter Latin had.}

videre, and the e that other verbs dropped but ver kept.
\end{chapteropening}

\section[ver]{ver — learn the verb before its past}

\label{lesson:ES-C16-ver}



\begin{quote}
\textbf{Ver} means “to see.” Its written \textbf{v} uses the Spanish sound you already know: it is not contrasted with a separate English-style \textbf{v} sound. Say \textbf{ver}.
\end{quote}

\begin{grammarlens}[title={Grammar Lens: three present forms}]
\noindent
\begin{tabularx}{\linewidth}{@{}>{\raggedright\arraybackslash}X>{\raggedright\arraybackslash}X@{}}
\toprule
\textbf{person} & \textbf{present} \\
\midrule
yo & \textbf{veo} \\
tú & \textbf{ves} \\
él / ella / usted & \textbf{ve} \\
\bottomrule
\end{tabularx}

The \textbf{tú} and third-person forms resemble the familiar \textbf{-er} row. The \textbf{yo} form keeps an extra vowel: \textbf{veo}. Reuse known \textbf{café}: \textbf{Veo café} — “I see coffee.”
\end{grammarlens}

\begin{cousinweb}
Spanish \textbf{ver} continues Latin \textbf{vidēre}, “to see,” through Old Spanish \textbf{veer}. Loss of the consonant between vowels helps explain the shorter shape; the two adjacent vowels later contracted in the infinitive. English \textbf{video} belongs to the same Latin family, a memory bridge rather than a new Spanish word for this lesson.
\end{cousinweb}

\subsection*{Your turn}
Say these aloud:

\begin{itemize}
\raggedright
  \item “ver — to see”
  \item “veo, ves, ve”
  \item “ver comes through veer from vidēre”
\end{itemize}

\begin{tcolorbox}[breakable,colback=teal!4,colframe=teal!35!black,title={Before you move on}]
What does \textbf{ver} mean? (“To see.”) Give its three learned present forms. (\emph{Veo, ves, ve.}) Which Latin verb stands behind it? (\textbf{Vidēre}.) Next: add \textbf{ser} to the open past.
\end{tcolorbox}
